%%%%%%%%%%%%%%%%%%%%%%%%%%%%%%%%%%%%%%%%%%%%%%%%%%%%%%%%%%%%%%%%%
% Presentacion 2/3: Estructuras de Datos en Python
% Serie "Python desde C" - TP4 Parte B (UNCuyo - Informatica y Programacion)
% Compilar con: pdflatex -shell-escape 02-estructuras-de-datos.tex  (dos pasadas)
%%%%%%%%%%%%%%%%%%%%%%%%%%%%%%%%%%%%%%%%%%%%%%%%%%%%%%%%%%%%%%%%%
\documentclass{beamer}

% Paquetes basicos
\usepackage[utf8]{inputenc}
\usepackage[T1]{fontenc}
\usepackage[spanish]{babel}

% Paquete para el resaltado de codigo de alta calidad
\usepackage{minted}

% Tema visual de la presentacion
\usetheme{Madrid}
\usecolortheme{beaver}
\setbeamertemplate{navigation symbols}{}

% Recapitula el indice, resaltando la seccion actual, al empezar cada seccion
\AtBeginSection[]{
  \begin{frame}{Indice de Contenidos}
    \tableofcontents[currentsection]
  \end{frame}
}

% --- Informacion del Titulo ---
\title[Python desde C -- 2/3]{Estructuras de Datos en Python}
\subtitle{2/3 -- Listas, tuplas y diccionarios (vs. arreglos y structs de C)}
\author{Informática y Programación}
\date{\today}
\institute{UNCuyo -- Facultad de Ingeniería}

% --- Inicio del Documento ---
\begin{document}

% --- Diapositiva del Titulo ---
\begin{frame}
  \titlepage
\end{frame}

% --- Diapositiva del Indice ---
\begin{frame}{Índice de Contenidos}
  \tableofcontents
\end{frame}

\begin{frame}{Dónde estamos}
  \begin{block}{Ya vimos (Presentación 1)}
    Sintaxis, tipos, condicionales, bucles, entrada de datos y funciones
    con parámetros simples.
  \end{block}
  \begin{alertblock}{Lo que falta para resolver el TP4}
    Con variables sueltas alcanza para poco: los tres pisos del parking,
    el tablero de la batalla naval, el inventario de la librería, todos
    necesitan agrupar datos. En C eso se resuelve con \texttt{arreglos} y
    \texttt{structs}. En Python: \textbf{listas, tuplas y diccionarios}.
  \end{alertblock}
\end{frame}

% ===============================================================
\section{1. Listas}
% ===============================================================

\begin{frame}[fragile]{Listas vs. arreglos de C}
  \begin{block}{Diferencia Principal: Tamaño}
    \begin{itemize}
      \item \textbf{Un arreglo de C tiene tamaño fijo}, decidido al declararlo (o al reservar memoria con \texttt{malloc}).
      \item \textbf{Una lista de Python crece y se achica sola}: \texttt{append}, \texttt{insert}, \texttt{remove} cambian su tamaño en tiempo de ejecución, sin \texttt{malloc}/\texttt{realloc}.
    \end{itemize}
  \end{block}
  \begin{columns}[T]
    \begin{column}{.48\textwidth}
      \centering\textbf{C}
      \begin{minted}[fontsize=\small, frame=lines, linenos]{c}
int notas[3] = {8, 5, 9};
int n = 3; // tamanio fijo
      \end{minted}
    \end{column}
    \begin{column}{.48\textwidth}
      \centering\textbf{Python}
      \begin{minted}[fontsize=\small, frame=lines, linenos]{python}
notas = [8, 5, 9]
notas.append(7)
# notas -> [8, 5, 9, 7]
      \end{minted}
    \end{column}
  \end{columns}
\end{frame}

\begin{frame}[fragile]{Indexado, \texttt{len()} y slicing}
  \begin{block}{\texttt{len()} en vez de llevar la cuenta a mano}
    En C solés guardar el tamaño en una variable aparte (\texttt{n}). En
    Python, \texttt{len(lista)} siempre está actualizado. El slicing
    (\texttt{lista[a:b]}) no tiene equivalente directo en C.
  \end{block}
  \begin{minted}[fontsize=\small, frame=lines, linenos]{python}
notas = [8, 5, 9, 7]

print(notas[0])     # primer elemento: 8
print(notas[-1])    # ultimo elemento: 7 (no existe en C)
print(len(notas))   # cantidad de elementos: 4
print(notas[1:3])   # sublista: [5, 9]
  \end{minted}
\end{frame}

\begin{frame}[fragile]{Recorrer una lista}
  \begin{block}{Dos formas, según si necesitás el índice}
    Si no necesitás la posición, \texttt{for elemento in lista} es más
    directo que el \texttt{for} de C. Si necesitás el índice, \texttt{range(len(...))} lo replica.
  \end{block}
  \begin{columns}[T]
    \begin{column}{.48\textwidth}
      \centering\textbf{C}
      \begin{minted}[fontsize=\small, frame=lines, linenos]{c}
int notas[3] = {8, 5, 9};
for (int i = 0; i < 3; i++) {
    printf("%d\n", notas[i]);
}
      \end{minted}
    \end{column}
    \begin{column}{.48\textwidth}
      \centering\textbf{Python}
      \begin{minted}[fontsize=\small, frame=lines, linenos]{python}
notas = [8, 5, 9]
for nota in notas:
    print(nota)

# con indice, si hace falta:
for i in range(len(notas)):
    print(i, notas[i])
      \end{minted}
    \end{column}
  \end{columns}
\end{frame}

\begin{frame}{Cuidado: aliasing con listas mutables}
  \begin{alertblock}{El bug más común al armar matrices en Python}
    Multiplicar una lista con \texttt{*} no copia sus elementos si son a
    su vez listas: \texttt{[[0]*4]*4} crea \textbf{4 referencias a la
    misma fila}, no 4 filas independientes. Modificar una, modifica las
    cuatro. No tiene equivalente en C: \texttt{int t[4][4]} siempre
    reserva memoria contigua real, fila por fila.
  \end{alertblock}
\end{frame}

\begin{frame}[fragile]{Aliasing: el ejemplo (tablero de la batalla naval)}
  \begin{minted}[fontsize=\small, frame=lines, linenos]{python}
# mal: las 4 filas son la MISMA lista, repetida por referencia
tablero = [[0] * 4] * 4
tablero[0][0] = 9
# tablero[1][0], [2][0] y [3][0] tambien cambiaron a 9!

# bien: una lista nueva por cada fila
tablero = [[0] * 4 for _ in range(4)]
tablero[0][0] = 9
# solo tablero[0][0] cambio
  \end{minted}
\end{frame}

% ===============================================================
\section{2. Tuplas}
% ===============================================================

\begin{frame}[fragile]{Tuplas: listas inmutables}
  \begin{block}{Se parecen a una lista, pero no se pueden modificar}
    Se escriben con paréntesis y, una vez creadas, no se les puede
    agregar, quitar ni reasignar elementos — como un arreglo de C
    declarado \texttt{const}, pero sin tamaño fijo previo.
  \end{block}
  \begin{minted}[fontsize=\small, frame=lines, linenos]{python}
punto = (3, 4)
print(punto[0])   # 3

punto[0] = 5      # TypeError: no se puede modificar una tupla
  \end{minted}
\end{frame}

\begin{frame}[fragile]{Retorno múltiple: el verdadero motivo para usarlas}
  \begin{alertblock}{En C, devolver más de un valor es incómodo}
    Una función en C solo devuelve \textbf{un} valor con \texttt{return}.
    Para devolver varios hay que usar punteros de salida (parámetros que
    la función modifica) o empaquetarlos en un \texttt{struct}.
  \end{alertblock}
  \begin{columns}[T]
    \begin{column}{.48\textwidth}
      \centering\textbf{C (punteros de salida)}
      \begin{minted}[fontsize=\footnotesize, frame=lines, linenos]{c}
void minmax(int v[], int n,
            int *min, int *max) {
    *min = *max = v[0];
    for (int i = 1; i < n; i++) {
        if (v[i] < *min) *min = v[i];
        if (v[i] > *max) *max = v[i];
    }
}
      \end{minted}
    \end{column}
    \begin{column}{.48\textwidth}
      \centering\textbf{Python (tupla)}
      \begin{minted}[fontsize=\footnotesize, frame=lines, linenos]{python}
def minmax(v):
    return min(v), max(v)

# desempaquetado directo:
lo, hi = minmax([8, 5, 9, 7])
# lo -> 5, hi -> 9
      \end{minted}
    \end{column}
  \end{columns}
\end{frame}

% ===============================================================
\section{3. Diccionarios}
% ===============================================================

\begin{frame}[fragile]{Diccionarios vs. structs de C}
  \begin{block}{Campos fijos vs. claves dinámicas}
    \begin{itemize}
      \item Un \texttt{struct} de C define sus campos \textbf{en tiempo
      de compilación} — no se pueden agregar campos nuevos mientras corre el programa.
      \item Un \texttt{dict} de Python asocia claves con valores \textbf{en
      tiempo de ejecución} — se pueden agregar, quitar o consultar claves libremente.
    \end{itemize}
  \end{block}
  \begin{columns}[T]
    \begin{column}{.48\textwidth}
      \centering\textbf{C}
      \begin{minted}[fontsize=\footnotesize, frame=lines, linenos]{c}
typedef struct {
    int codigo;
    char nombre[50];
    float precio;
    int stock;
} Libro;

Libro l = {1, "Rayuela",
           15000.0, 10};
      \end{minted}
    \end{column}
    \begin{column}{.48\textwidth}
      \centering\textbf{Python}
      \begin{minted}[fontsize=\footnotesize, frame=lines, linenos]{python}
libro = {
    "codigo": 1,
    "nombre": "Rayuela",
    "autor": "Cortazar",
    "precio": 15000.0,
    "stock": 10,
}
      \end{minted}
    \end{column}
  \end{columns}
\end{frame}

\begin{frame}[fragile]{Usar y recorrer un diccionario}
  \begin{block}{Acceso por clave, no por índice}
    Se lee y escribe con \texttt{d["clave"]}, no con una posición numérica.
    Para recorrer todo el contenido, \texttt{.items()} da clave y valor juntos.
  \end{block}
  \begin{minted}[fontsize=\small, frame=lines, linenos]{python}
libro = {"codigo": 1, "nombre": "Rayuela", "stock": 10}

print(libro["nombre"])     # Rayuela
libro["stock"] -= 2        # vender 2 unidades

for clave, valor in libro.items():
    print(clave, "->", valor)
  \end{minted}
\end{frame}

\begin{frame}[fragile]{Inventario completo: lista de diccionarios}
  \begin{block}{El patrón para el ejercicio 3 del TP4}
    Un \texttt{struct Libro[]} en C es una lista de diccionarios en
    Python: cada libro es un \texttt{dict}, y el inventario entero es una
    \texttt{list} de esos \texttt{dict}.
  \end{block}
  \begin{minted}[fontsize=\small, frame=lines, linenos]{python}
inventario = [
    {"codigo": 1, "nombre": "Rayuela", "stock": 10},
    {"codigo": 2, "nombre": "El Aleph", "stock": 3},
]

for libro in inventario:
    print(libro["codigo"], libro["nombre"], libro["stock"])
  \end{minted}
\end{frame}

% ===============================================================
\section{4. Resumen y Conclusiones}
% ===============================================================

\begin{frame}{Tabla Comparativa Rápida}
  \small
  \begin{tabular}{|p{2.6cm}|p{3.4cm}|p{3.4cm}|}
    \hline
    \textbf{Necesito...} & \textbf{Lenguaje C} & \textbf{Lenguaje Python} \\
    \hline \hline
    Colección de tamaño fijo & \texttt{int v[N]} & lista (aunque no cambie de tamaño) \\
    \hline
    Colección que crece & \texttt{malloc}/\texttt{realloc} & lista con \texttt{append} \\
    \hline
    Colección inmutable & \texttt{const int v[N]} & tupla \\
    \hline
    Registro con campos fijos & \texttt{struct} & diccionario \\
    \hline
    Devolver varios valores & Punteros de salida o \texttt{struct} & tupla (\texttt{return a, b}) \\
    \hline
  \end{tabular}
\end{frame}

\begin{frame}{Lo que sigue}
  \begin{block}{Con esto ya podés}
    Modelar los pisos del parking (listas), el tablero de la batalla naval
    (listas anidadas), y el inventario de la librería (lista de
    diccionarios) — la base de los ejercicios 5, 6 y 3 del TP4.
  \end{block}
  \begin{alertblock}{Presentación 3: Cadenas, Archivos y Funciones Avanzadas}
    Falta lo que necesitan Wordle (comparar letra por letra), el Juego de
    la Vida (recorrer una grilla) y las estadísticas del archivo de texto
    (leer un \texttt{.txt} y contar palabras): \textbf{manejo avanzado de
    cadenas, lectura de archivos y funciones con retorno de colecciones}.
    Eso es lo que viene.
  \end{alertblock}
\end{frame}

\begin{frame}
  \vfill
  \centering
  \begin{beamercolorbox}[sep=8pt,center,shadow=true,rounded=true]{title}
    \Huge \bfseries ¡A programar en Python!
  \end{beamercolorbox}
  \vfill
\end{frame}

\end{document}
